% SEGMENT OLP-0001-S01
\chapter*{Open Logic Project பற்றிய அறிமுகம்}
\addcontentsline{toc}{chapter}{Open Logic Project பற்றிய அறிமுகம்}

\textit{Open Logic Text} என்பது, முறைசார்ந்த மேல்தருக்கவியலையும்
முறைசார்ந்த வழிமுறைகளையும் பற்றிய, திறந்த மூலமாகக் கூட்டாக
உருவாக்கப்படும் பாடநூல். இது இடைநிலையிலிருந்து தொடங்குகிறது
(அதாவது, முறைசார்ந்த தருக்கத்திற்கான அறிமுகப் பாடத்திற்குப் பிறகு).
கணிதத் துறையைச் சாராத வாசகர்களை, குறிப்பாக மெய்யியல் மற்றும்
கணினி அறிவியல் மாணவர்களை, நோக்கமாகக் கொண்டிருந்தாலும்,
இந்நூல் துல்லியமான நிறுவல் முறையைக் கடைப்பிடிக்கிறது.

% SEGMENT OLP-0001-S02
தற்போது சேர்க்கப்பட்டுள்ள சில தலைப்புகளின் பாடப்பொருள் இன்னும்
முழுமை பெறாமல் இருக்கலாம்; பல பிரிவுகளுக்கு மேலும் விரிவான
திருத்தம் தேவைப்படுகிறது. எதிர்காலத்தில் மேலும் பல தலைப்புகளைச்
சேர்க்கத் திட்டமிடுகிறோம். கலைச்சொல் அகராதி, மேலும் வாசிப்பதற்கான
நூற்பட்டியல், வரலாற்றுக் குறிப்புகள், படங்கள், மேம்பட்ட விளக்கங்கள்,
முடிவுகள் மெய்யியல், கணினி அறிவியல் மற்றும் கணிதத்திற்கு எவ்வாறு
பொருந்துகின்றன என்பதை விளக்கும் பிரிவுகள், மேலும் பயிற்சிகளும்
எடுத்துக்காட்டுகளும் போன்றவற்றைச் சேர்க்கவும் திட்டமிடுகிறோம்.
பிழை ஒன்றைக் கண்டாலோ பரிந்துரை ஒன்றைத் தர விரும்பினாலோ,
\href{https://github.com/OpenLogicProject/OpenLogic/wiki/Contributing}{திட்டக் குழுவிற்குத் தெரிவிக்கவும்}.

% SEGMENT OLP-0001-S03
இத்திட்டம் திறந்த மூலக் கொள்கையின் அடிப்படையில் செயல்படுகிறது.
நூல் இலவசமாகக் கிடைப்பதுடன், அதன் LaTeX மூலத்தையும்
Creative Commons Attribution உரிமத்தின்கீழ் வழங்குகிறோம்.
உரிய படைப்பாளர் குறிப்பைத் தரும் நிபந்தனையுடன், எவரும்
எங்கள் படைப்பைப் பதிவிறக்கவும், பயன்படுத்தவும், மாற்றவும்,
மறுவரிசைப்படுத்தவும், வேறு வடிவத்திற்கு மாற்றவும், மறுபகிர்வு
செய்யவும் அந்த உரிமம் அனுமதிக்கிறது. மேலும் விவரங்களுக்கு,
Open Logic Project இணையதளமான
\href{http://openlogicproject.org/}{openlogicproject.org} ஐப் பார்க்கவும்.
